% OWNER: sections/s4_1_tests.tex. Edit only this file.
\seclabel
\claim{\secnum{4.1}{sec:wrappertests}\quad Start-up, the feed, a buy and a sell worked; \code{plan()} needed a 100\% cap}

\lead{The wrapper, built to the design of Section~\ref{sec:design}, ran against the live server from 09:27 to 09:55 KST on 24 Nov 2026, with the floor test of Section~\ref{sec:oneshare} at 10:26 and 10:33.} The package \code{pytimefolio} holds the modules of \figref{modules}. A strategy uses one class, \code{Timefolio}, from the \code{client} module (\figref{client}). Price arguments (\code{**price}) are \code{limit_prc}, \code{limit_idx} or \code{stop_prc}, exactly one of them, with optional \code{hm0} and \code{hm1} (\figref{orderapi}).

% OWNER: figs/client.tex. Edit only this file.
\def\fclL#1#2{\hspace*{#1mm}\code{#2}}% one code line, indented #1 mm
\def\fclC#1{\leavevmode{\rmfamily\color{greyD}#1}}% comment beside a code line
\fig[wide]{client}{One class gives the strategy its state, its orders and the guideline status}[Use of the \code{Timefolio} client, version 0.1.0; comments on the right]{%
{\fontsize{7.6}{9.6}\selectfont\begin{tabular}{@{}>{\RaggedRight\arraybackslash}p{106mm} >{\RaggedRight\arraybackslash}p{70mm}@{}}
\fclL{0}{async with Timefolio(91, email=os.environ["EMAIL"],} & \fclC{competition 91; the password is passed in,}\\
\fclL{28.4}{password=os.environ["PASSWORD"]) as tf:} & \fclC{never stored; the cookie jar is saved}\\[2pt]
\fclL{4}{w = tf.weights()} & \fclC{\code{prodId} $\rightarrow$ current \code{wei} (\%)}\\
\fclL{4}{tf.orders_left()} & \fclC{80 minus today's new orders and stops}\\[2pt]
\fclL{4}{orders = tf.plan(targets, min_wei=0.05, max_total=100,} & \fclC{dry run: sends nothing}\\
\fclL{27}{prices={"A000660": p}, limit_idx=5)} & \fclC{\code{ValueError} on a rule breach}\\
\fclL{4}{tf.dropped} & \fclC{orders below one share, left out}\\[2pt]
\fclL{4}{try:} & \\
\fclL{8}{ids = await tf.rebalance(targets, limit_idx=5)} & \fclC{\code{plan()}, then the orders 1\,s apart}\\
\fclL{4}{except RebalanceError as e:} & \\
\fclL{8}{e.placed} & \fclC{order IDs sent before the failure}\\[2pt]
\fclL{4}{i = await tf.buy("A005930", 0.1, limit_prc=150700)} & \fclC{or \code{tf.sell}; returns the order ID}\\
\fclL{4}{await tf.stop(i)} & \fclC{\code{stop(*ids)} takes several}\\
\fclL{4}{r = await tf.wait(i)} & \fclC{\code{r.outcome}, \code{r.cum_qty}, \code{r.qty2ord}, \code{r.avg_px}}\\[2pt]
\fclL{4}{g = await tf.guidelines()} & \fclC{\code{g["ok"]}, \code{g["sector"]}, \code{g["turnover"]}}\\
\end{tabular}}}{\code{targets} maps tickers to target weights in percent and describes the whole portfolio: a held stock left out of it gets target 0 and is sold. \code{plan()} returns the orders \code{rebalance()} would send, sells first, after subtracting what working orders still have to fill; \code{prices} gives the price of a stock the portfolio does not hold, so that its one-share check can run (Section~\ref{sec:oneshare}). Price arguments (\code{**price}) are exactly one of \code{limit_prc}, \code{limit_idx} or \code{stop_prc}, with optional \code{hm0} and \code{hm1}. \code{buy()} and \code{sell()} return the order ID found in the feed, since \code{AddOrder} returns none (Section~\ref{sec:orders}). \code{rebalance()} first re-counts today's orders on the server and sends nothing if its orders exceed what is left of the daily cap. \code{r.outcome} is \code{filled}, \code{stopped} or \code{partly_filled_then_stopped}. \code{p} stands for a price in KRW supplied by the caller.}{DART Quant Team; \code{client.py} of pytimefolio 0.1.0}


\lead{Entering the \code{async with} block runs 5 steps before the strategy gets control.} First, the client renews the access token from the cookie jar stored in \code{session.json} and logs in only if that renewal fails (D17). Second, it chooses the portfolio by the rule of Section~\ref{sec:queries}. Third, it starts the live feed as a background task and waits up to 30\,s for the first snapshot and fails start-up with a timeout if none arrives, so \code{weights()} never reads an empty state; afterwards every call that reads the state or places an order raises \code{StaleFeed} when the last feed reply is older than 26\,s (2 $\times$ 3\,s + 20\,s), so stale weights are never used. Fourth, it loads \code{Universe} once to learn each stock's sector (Section~\ref{sec:guidelines}). Fifth, it counts the day's orders in \code{Portfolio/Orders}, new orders plus stops, and seeds the daily cap with the larger of that number and its own count file, so a restart during the day does not reset the count to zero. Leaving the block cancels the feed task and closes the WebSocket and the HTTP client without calling \code{Logout} (D16).

\Needspace{3\baselineskip}\lead{Each module passed its live test, and 31 offline tests cover the arithmetic and the failure paths that a live server cannot be made to produce (\figref{tests}).} Nine of the offline tests check pure functions; the other 22 replace the server with \code{httpx.MockTransport} and fakes and produce what the live server would not on demand: each error status, a lost \code{AddOrder} reply, a resend that fails too, a feed that sends only pings, a stale feed, 3 concurrent writes (1.002\,s apart at the mock server), 2 identical orders in a row, 2 new records that match one order, and a 403 or 429 on a read or a write followed by a 200. The note to \figref{tests} lists them. All 31 run in GitHub Actions on every push and pull request, on Python 3.11 to 3.13; the live checks in \code{tests/live} never run in CI.

% OWNER: figs/tests.tex. Edit only this file.
\fig[wide]{tests}{Every test passed after the 2 defects the live tests exposed were fixed}[Wrapper tests by goal (\figref{goals}): live against the competition server, 09:27 to 10:34 KST on 24 Nov 2026, one portfolio, and offline]{%
{\tabfont\begin{tabularx}{\linewidth}{>{\RaggedRight\arraybackslash}p{16mm} >{\RaggedRight\arraybackslash}p{25mm} >{\RaggedRight\arraybackslash}p{47mm} >{\RaggedRight\arraybackslash}X}
\thd{Goal} & \thd{Test} & \thd{What was checked} & \thd{Result}\\
G1 & \code{transport} & Cost of the first request, of a second client's first request and of a reused connection; clock offset from the \code{Date} header & First request of a fresh process 1.2 to 1.3\,s; a second client's first request 130\ms; reused 31 to 39\ms. Server clock minus PC clock $-2.98$\,s, matching the 3\,s of Section~\ref{sec:auth}\\\rowsep
G1 & \code{auth} & Login, cookie storage and renewal across 2 runs & First run logged in and saved the cookie jar to \code{~/.config/pytimefolio/session.json} with mode 0600; second run renewed the token without logging in\\\rowsep
G2, G5 & \code{rest} & Portfolio rule; concurrent reads on one connection; ticker check & The live portfolio selected by the rule. \code{Session}, \code{Universe} (2,648), \code{Trades} (64) and \code{Orders} (5 that morning) fetched concurrently on one HTTP/2 connection. \code{ValidateProduct} for \code{A005930} returned no warnings\\\rowsep
G2 & \code{hub}, \code{state} & One connection polled every 3\,s; delta merge & 1 snapshot, then 20 deltas without a drop. Merged state equalled same-\code{t} snapshots in 15 of 15 rounds (the test of Section~\ref{sec:feed})\\\rowsep
G3 & \code{orders}, buy & Samsung Electronics 0.1\%, limit KRW~150,700, 20\% below the market price of KRW~188,400 & Found in the feed 1 to 3\,s after \code{AddOrder}: Id 1650412, \code{Working}, \code{qty2ord} 4. Stopped after 3\,s; \code{wait()} returned stopped, 0 of 4\\\rowsep
G3 & \code{orders}, sell & LIG Nex1 0.1\%, limit KRW~427,000, 20\% above the market price of KRW~355,500 & Order record \code{ls} = L, \code{ex} = X, \code{sgn} = $-1$, \code{qty2ord} 2. Stopped; result stopped, 0 of 2\\\rowsep
\hl G3 & \code{plan()}, total & Samsung Electronics 50\% against its \code{weiLimit} of 40\% & Capped at 40\%; water-filling spread the excess over the other 22 holdings. The targets summed to 139\% and \code{plan()} accepted them: defect 1\\\rowsep
G3 & \code{plan()}, one share (\ref{sec:oneshare}) & 3 small orders with \code{min_wei} 0.01, 09:52 & SK hynix +0.03\% dropped with a warning (one share 0.1413\%); Samsung Electronics +0.1\% and Leeno Industrial +0.06\% kept\\\rowsep
\hl G3 & \code{guidelines()} (\ref{sec:guidelines}) & Feed fields and the \code{Violations} reply, 09:52 & IT shown over its sector limit while the server counted 0 violations: defect 2. After the fix, 0 server violations; \code{ok} False on turnover 4.8\%\\\rowsep
G3, G5 & \code{plan()}, sector (\ref{sec:guidelines}) & 3 target sets, 09:55 & Unchanged weights: 0 orders. Industrials +6\,pp: rejected. Industrials +3\,pp: 2 orders (\figref{sectortest})\\\rowsep
G2 & \code{live_fields} & \code{bs_updated} in each delta for 60\,s, 10:33; no orders & \code{bs_updated} arrived as a dict in 20 of 20 deltas; \code{na} moved from KRW~844,505,167 to 844,076,987, so net assets stay current\\\rowsep
G3 & \code{live_rounding} (\ref{sec:oneshare}) & Samsung Electronics 0.04\%, worth 1.76 shares, limit 20\% below the market; stopped at once & \code{qty2ord} 1 both times (Id 1650693 at about 10:26, Id 1650719 at 10:33:45), 0 filled: the server rounds down\\\rowsep
G3 & Corrected code, live & Orders placed and stopped by the code changed after the review: sizes quantised to 0.01\%; an order the server may have received never resent & Order records store \code{wei} as sent (0.04, 0.1 and 0.3 seen). Order 1650719 was found in the feed as \code{Accepted}, stopped at once and ended \code{Done}. Counter at 10:34: 9 new orders and 7 stops, 16 of 80 used\\\rowsep
G1 to G3, G5 & Offline & 31 unit tests (9 pure, 22 failure paths with mocked HTTP) & All pass\\
\end{tabularx}}}{Market prices are those of the feed at the time of the order. Test orders had limits 20\% from the market so that they could not fill. Highlighted rows are the 2 defects. Defect 1: water-filling keeps the total, so a 139\% target stayed 139\% after capping; \code{plan()} now raises \code{ValueError} when targets sum above \code{max_total}, 100\% by default. Defect 2: sector limits were keyed by the \code{msecs} record \code{Id} instead of the sector code \code{sec}, so IT was compared with another sector's limit; keying on \code{sec} fixed it. The tests of \ref{sec:oneshare} at 09:52 and of \ref{sec:guidelines} called \code{plan()} and \code{guidelines()} only and sent no orders. Offline tests, pure: state merge (string keys converted to \code{int}, unchanged fields kept, null reply, added and removed records, snapshot replacement), water-filling, target-to-order conversion, order body against the request captured at 09:18, \code{sector_violations}, KRX tick sizes. Failure paths (\code{httpx.MockTransport} and fakes): status codes 401, 403, 429, 400 and 502 mapped to exceptions; the transport retries a GET once on a connection failure and never a POST, and \code{AddOrder} is sent once more only when the connection was never established; a read error counted as sent; 5 concurrent \code{token()} calls cause 1 \code{SilentRefresh}; a 401 renews once; a corrupt session file leads to a login; size quantisation; working orders not ordered again; daily cap and sync; a connection failure not counted; 3 concurrent writes reach the server at least 1\,s apart (1.002\,s measured); a 403 and a 429 each halt all writes, even when the server answers 200 again; a lost \code{AddOrder} reply not resent; a stale feed blocks orders and \code{wait()}; feed freshness. Added after the live tests: the cookie expiry kept across a restart; a failed resend still searched, with no third send; the same order placed twice gets 2 IDs; 2 new records matching one order raise \code{AmbiguousState}; a 403 on a read halts writes; server pings do not extend the reply deadline; a blocked feed stops reconnecting. They run in CI on Python 3.11 to 3.13; \code{tests/live} does not.}{DART Quant Team; \code{tests/} and \code{tests/live/} of pytimefolio 0.1.0}


\lead{A code review before these tests found 7 defects in the order and feed paths, and all 7 are fixed in the code the tests ran.} The fixes are the behaviour already described: (1) \code{wei} quantised once to 0.01\% and matched as sent; (2) no automatic resend of an order whose reply was lost, with one lock across write and search (Section~\ref{sec:orders}); (3) \code{StaleFeed} after 26\,s; (4) reconnection with full-jitter backoff; (5) the merge of \code{bs_updated} (Section~\ref{sec:feed}); (6) the subtraction of working orders in \code{plan()} (Section~\ref{sec:guidelines}); and (7) a daily count that includes stops and survives a restart, with the halt of all writes after a 403 or 429 (Section~\ref{sec:design}).

\lead{The sell test settles how a sell is sent.} The wrapper's \code{sell()} sends the same \code{AddOrder} body as a buy with the exit code \code{ex} = X in place of E, and \code{ls} stays L. The server accepted it and recorded the sign \code{sgn} = −1 (a sell) with \code{qty2ord} 2, so a sell is sent as an exit (\code{ex} = X) with the same side code (\code{ls} = L) as a buy. Section~\ref{sec:orders}'s orders were all buys; with this test both sides have been placed, found in the feed and stopped.

\lead{The first \code{plan()} test found that a target set summing above 100\% passed through.} Water-filling moves the weight cut at a limit onto the other stocks and keeps the total unchanged, so a 139\% target stayed 139\% after capping and \code{plan()} would have asked for a portfolio worth 139\% of its value. \code{plan()} now raises \code{ValueError} before capping when the targets sum above \code{max_total}, 100\% by default (D19).

\lead{By 10:34 the day's count stood at 9 new orders and 7 stops, and the wrapper's counter showed 16 of its 80 used.} Section~\ref{sec:orders}'s tests used 5 orders and 3 stops, the buy and the sell added 2 of each, and the floor test of Section~\ref{sec:oneshare} added 2 more of each; the tests of \code{plan()} and \code{guidelines()} sent none. Until the organiser says whether cancels count toward the 100 (question Q2 of \figref{openq}, Section~\ref{sec:validity}), the wrapper counts a stop as an order, so 64 orders remain under its cap and the organiser's threshold of 100 is 84 away.

\lead{The tests show G2 and G5; G1 and G3 are partly shown and G4 is still open (\figref{goalstatus}).} G1 needs the wrapper to stay logged in until 30 Nov, and the tests covered one morning: the second run renewed from the stored cookie jar without logging in, but no process has yet run past 720\,s and renewed inside the same run, and none has run for 7 days. G2 is judged in 2 parts: the merged state equalled same-\code{t} snapshots in 15 of 15 rounds (Section~\ref{sec:feed}), observed live; the refusal of a state older than 26\,s is shown offline only, since no live outage was produced. G3 is partly shown because the 15\,s confirmation window rests on live buys found in 1 to 3\,s and was not tested under a slow server, and because \code{plan()} checks target weights at planning time, not weights during execution. G4, that polling never delays an order, has no outcome measure yet; the tests show only that the feed and the orders share one event loop. G5 has passing live and offline tests.

% OWNER: figs/goalstatus.tex. Edit only this file.
\fig[wide]{goalstatus}{The tests show G2 and G5; G1 and G3 are partly shown and G4 is not measured}[Each goal of \figref{goals} against the tests of \figref{tests} that bear on it, 24 Nov 2026; G2's two parts are judged separately]{%
{\tabfont\begin{tabularx}{\linewidth}{>{\RaggedRight\arraybackslash}p{34mm} >{\RaggedRight\arraybackslash}p{46mm} >{\RaggedRight\arraybackslash}X >{\RaggedRight\arraybackslash}p{15mm}}
\thd{Goal} & \thd{Test} & \thd{Result} & \thd{Status}\\
\hl G1 Stay logged in for the whole competition without a person & \code{auth} live across 2 processes; offline: 5 concurrent \code{token()} calls, a 401 retry, a corrupt session file, the cookie expiry across a restart & The second process renewed from the stored cookie jar without logging in; offline, 5 calls caused 1 \code{SilentRefresh} and the expiry survived a restart. Renewal after 720\,s inside one process and the login before the 7-day cookie expiry have not been exercised & Partly shown\\\rowsep
G2(a) Merge: merged state equals the same-\code{t} snapshot & \code{hub}, \code{state} and \code{live_fields} live & Merged \code{pos} and \code{ords} equalled same-\code{t} snapshots in 15 of 15 rounds; \code{bs_updated} in 20 of 20 deltas kept \code{na} current. \code{bs} and \code{msecs} not compared field by field & Shown\\\rowsep
G2(b) Freshness: no read or order from a state older than 26\,s & Offline: stale-feed test, feed freshness, pings that do not extend the reply deadline & A feed whose last reply is older than 26\,s raises \code{StaleFeed} and blocks orders and \code{wait()}; a feed that only pings times out after the reply deadline. No live outage was produced & Shown offline\\\rowsep
\hl G3 Place, stop and track orders within the organiser's limits and the competition's guidelines & Offline: daily cap and sync, write spacing, never resend, failed resend, 2 identical orders, halt on 403 and 429; live: buy, sell, \code{plan()} and the corrected code & Writes past the cap refused; a lost \code{AddOrder} reply not resent; 3 concurrent writes reached the mock server 1.002\,s apart; all writes halted after a 403 or 429; live orders placed, found in 1 to 3\,s and stopped, 16 of 80 used by 10:34. Not shown: confirmation within 15\,s under a slow server (live finds were buys only), and limits during execution, since \code{plan()} checks target weights at planning time only & Partly shown\\\rowsep
\hl G4 Polling never delays an order & None & Feed polling and orders share one event loop; the delay of an order while a poll is in flight has not been measured & Not shown\\\rowsep
G5 Stop on ambiguity instead of guessing & \code{rest} live (portfolio rule); offline lost-reply, failed-resend and status tests; \code{plan()} sector and total tests & Portfolio chosen by the rule; an order not found raises \code{AmbiguousState} without a resend; a 400 raises \code{RequestError}; \code{plan()} rejected the target set over the Industrials limit and, after defect 1, targets above 100\% & Shown\\
\end{tabularx}}}{Highlighted rows are the goals not fully shown. Shown: the tests listed met the goal's measure of success in \figref{goals}. Shown offline: the code's behaviour was shown against a mocked server but not observed live. Partly shown: part of the measure was met and the rest not yet exercised. Not shown: no test measured the goal. What each result lets a user rely on, and under which assumption, is in \figref{guarantees}. Live tests ran on one account on one morning (Section~\ref{sec:validity}).}{DART Quant Team; \figref{tests}}


\verdict{the wrapper's \code{plan()} rejects targets that sum above 100\% (D19). \code{sell()} sends \code{ls} = L with \code{ex} = X, as verified live (D21).}
